\documentclass{article}
\usepackage{amsmath,amssymb,amsthm}
\usepackage{algorithm}
\usepackage{algpseudocode}
\usepackage{geometry}
\geometry{margin=1in}
\newtheorem{theorem}{Theorem}
\newtheorem{lemma}{Lemma}
\newtheorem{assumption}{Assumption}
\newtheorem{corollary}{Corollary}
\newcommand{\argmin}{\operatorname*{arg\,min}}
\newcommand{\Breg}[2]{D_{h}(#1,#2)}
\newcommand{\inner}[2]{\langle #1,#2\rangle}
\newcommand{\grad}{\nabla}
\begin{document}

\section*{Algorithm Name}
\textbf{Mirror Descent (MD)} – a proximal‑gradient method defined with respect to a
strongly convex distance generating function $h$.

\vspace{0.5em}
\section*{Mathematical Formulation}
Let $f:\mathbb{R}^{d}\rightarrow\mathbb{R}$ be a convex differentiable objective,
$h:\mathbb{R}^{d}\rightarrow\mathbb{R}$ a $1$‑strongly convex prox‑function w.r.t.\ a norm
$\|\cdot\|$, and let the associated Bregman divergence be
\[
\Breg{x}{y}=h(x)-h(y)-\inner{\grad h(y)}{x-y}.
\]
Given a stepsize sequence $\{\eta_k\}_{k\ge 0}$, Mirror Descent performs for $k=0,1,\dots$
\[
\boxed{
x_{k+1}= \argmin_{x\in\mathbb{R}^{d}}
\Big\{\,\inner{\grad f(x_k)}{x}+ \frac{1}{\eta_k}\,\Breg{x}{x_k}\,\Big\}.
}
\tag{MD}
\]

The optimality condition of the above minimisation yields the equivalent
\[
\grad h(x_{k+1}) = \grad h(x_k)-\eta_k\,\grad f(x_k). \tag{1}
\]

\vspace{0.5em}
\section*{Pseudocode}
\begin{algorithm}[H]
\caption{Mirror Descent}
\begin{algorithmic}[1]
\Require Convex $f$, prox‑function $h$, stepsizes $\{\eta_k\}$, initial point $x_0$
\State Compute $u_0 \gets \grad h(x_0)$
\For{$k=0,1,\dots,T-1$}
    \State Compute gradient $g_k \gets \grad f(x_k)$
    \State Update dual variable $u_{k+1} \gets u_k - \eta_k\, g_k$
    \State Map back to primal: $x_{k+1} \gets (\grad h)^{-1}(u_{k+1})$
\EndFor
\State \textbf{return} $x_T$
\end{algorithmic}
\end{algorithm}

\vspace{0.5em}
\section*{Dry Run Example}
Consider the scalar problem 
\[
f(w) = (w-3)^2,\qquad h(w)=\tfrac12 w^{2} \;( \text{so } \Breg{w}{w'}=\tfrac12 (w-w')^{2}).
\]
The gradient $\grad f(w)=2(w-3)$ and $\grad h(w)=w$.
Take $w_0=0$ and a constant stepsize $\eta=0.1$.

\begin{center}
\begin{tabular}{c|c|c|c}
$k$ & $g_k=\grad f(w_k)$ & $w_{k+1}=w_k-\eta g_k$ & $f(w_{k+1})$\\ \hline
0 & $2(0-3)=-6$ & $0-0.1(-6)=0.6$ & $(0.6-3)^2=5.76$\\
1 & $2(0.6-3)=-4.8$ & $0.6-0.1(-4.8)=1.08$ & $(1.08-3)^2=3.6864$\\
2 & $2(1.08-3)=-3.84$ & $1.08-0.1(-3.84)=1.464$ & $(1.464-3)^2=2.3689$
\end{tabular}
\end{center}
The objective value decreases monotonically, illustrating a single‑step
Mirror Descent update (which coincides with ordinary gradient descent for this
quadratic $h$).

\vspace{1em}
\section*{Assumptions}
\begin{assumption}[Convexity of $f$]
$f$ is convex:
\[
f(y)\ge f(x)+\inner{\grad f(x)}{y-x},\qquad\forall x,y\in\mathbb{R}^{d}.
\tag{A1}
\]
\end{assumption}
\begin{assumption}[Smoothness w.r.t. $h$]
There exists $L>0$ such that for all $x,y$,
\[
f(y)\le f(x)+\inner{\grad f(x)}{y-x}+ \frac{L}{2}\,\Breg{y}{x}.
\tag{A2}
\]
\end{assumption}
\begin{assumption}[Strong convexity of $h$]
$h$ is $1$‑strongly convex w.r.t.\ $\|\cdot\|$:
\[
\Breg{x}{y}\ge \tfrac12\|x-y\|^{2},\qquad\forall x,y.
\tag{A3}
\]
\end{assumption}
\begin{assumption}[Stepsize]
We fix a constant stepsize $\eta_k\equiv\eta\le \frac{1}{L}$.
\tag{A4}
\end{assumption}

\vspace{1em}
\section*{Main Convergence Theorem}
\begin{theorem}[Deterministic Mirror Descent, $O(1/T)$ rate]\label{thm:md}
Let $x^{\star}\in\arg\min f$ and suppose Assumptions (A1)–(A4) hold.
Define the averaged iterate $\bar{x}_T = \frac1T\sum_{k=0}^{T-1}x_k$.
Then for any $T\ge1$,
\[
f(\bar{x}_T)-f(x^{\star})\;\le\;
\frac{\Breg{x^{\star}}{x_0}}{\eta T}.
\tag{2}
\]
In particular, with $\eta=1/L$,
\[
f(\bar{x}_T)-f(x^{\star})\;\le\;
\frac{L\,\Breg{x^{\star}}{x_0}}{T}=O\!\left(\frac{1}{T}\right).
\]
\end{theorem}

\vspace{1em}
\section*{Theoretical Properties}
\begin{itemize}
\item \textbf{Stability.} The update (1) is a non‑expansive map in the dual space
because $\grad h$ is $1$‑Lipschitz (by strong convexity of $h$). Hence
$\|u_{k+1}-u^{\star}\|_{*}\le\|u_k-u^{\star}\|_{*}+ \eta\|g_k\|_{*}$,
preventing divergence for any $\eta\le1/L$.
\item \textbf{Hyperparameter Sensitivity.} The bound (2) is monotone decreasing
in $\eta$ up to the limit $\eta=1/L$. Larger $\eta$ yields a tighter
$1/(\eta T)$ factor but violates the smoothness condition, breaking the
inequality in Lemma~\ref{lem:smooth}.
\item \textbf{Comparison to Gradient Descent.} For $h(w)=\frac12\|w\|^{2}$,
$\Breg{x}{y}=\frac12\|x-y\|^{2}$ and Mirror Descent reduces to (projected) gradient
descent with the same $O(1/T)$ rate. For non‑Euclidean $h$ (e.g., entropy),
MD adapts to geometry and can obtain smaller Bregman diameters,
thus a better constant in (2) compared with Euclidean GD.
\end{itemize}

\vspace{1em}
\section*{Discussion}
The $O(1/T)$ guarantee matches the optimal rate for deterministic first‑order
methods on smooth convex problems. The proof hinges on the three‑point
identity of Bregman divergences and the smoothness assumption (A2). Extensions
to stochastic gradients replace the exact gradient by an unbiased estimator and
add a variance term, leading to an $O(1/\sqrt{T})$ rate; the deterministic proof
remains the core of that analysis.

\vspace{1em}
\section*{Preliminaries (Definitions and Lemmas)}
\begin{lemma}[Three‑point identity]\label{lem:3pt}
For any $x,y,z$,
\[
\Breg{x}{z}= \Breg{x}{y}+ \Breg{y}{z}+ \inner{\grad h(y)-\grad h(z)}{x-y}.
\tag{3}
\]
\end{lemma}
\begin{proof}
Direct expansion of the definition of $\Breg{\cdot}{\cdot}$ yields (3). \hfill\qedsymbol
\end{proof}
\begin{lemma}[Smoothness inequality]\label{lem:smooth}
If (A2) holds then for any $x,y$,
\[
f(y)-f(x)-\inner{\grad f(x)}{y-x}\le \frac{L}{2}\,\Breg{y}{x}.
\tag{4}
\]
\end{lemma}
\begin{proof}
This is precisely the statement of (A2). \hfill\qedsymbol
\end{proof}
\begin{lemma}[Strong convexity of $h$]\label{lem:strong}
If (A3) holds, then for any $x,y$,
\[
\inner{\grad h(x)-\grad h(y)}{x-y}\ge \|x-y\|^{2}\ge 2\,\Breg{x}{y}.
\tag{5}
\]
\end{lemma}
\begin{proof}
From (A3), $\Breg{x}{y}\ge \frac12\|x-y\|^{2}$. Rearranging gives (5). \hfill\qedsymbol
\end{proof}

\vspace{1em}
\section*{Main Proof (Step‑by‑Step Derivation)}
We prove Theorem~\ref{thm:md}. Throughout the proof let $x^{\star}\in\arg\min f$.
All expectations are omitted because the algorithm is deterministic.

\bigskip
\noindent\textbf{[Step 1]}  Start from the optimality condition (1):
\[
\grad h(x_{k+1}) = \grad h(x_k)-\eta\,\grad f(x_k).
\tag{6}
\]

\noindent\textbf{[Step 2]}  Take the inner product of both sides of (6) with
$x_{k+1}-x^{\star}$:
\[
\inner{\grad h(x_{k+1})-\grad h(x_k)}{x_{k+1}-x^{\star}}
= -\eta\,\inner{\grad f(x_k)}{x_{k+1}-x^{\star}}.
\tag{7}
\]

\noindent\textbf{[Step 3]}  Apply the three‑point identity (Lemma~\ref{lem:3pt})
with $(x,y,z)=(x^{\star},x_{k+1},x_k)$:
\[
\Breg{x^{\star}}{x_k}= \Breg{x^{\star}}{x_{k+1}}+\Breg{x_{k+1}}{x_k}
+ \inner{\grad h(x_{k+1})-\grad h(x_k)}{x^{\star}-x_{k+1}}.
\tag{8}
\]
Rearrange (8) to isolate the inner product term:
\[
\inner{\grad h(x_{k+1})-\grad h(x_k)}{x_{k+1}-x^{\star}}
= \Breg{x^{\star}}{x_{k+1}}-\Breg{x^{\star}}{x_k}
- \Breg{x_{k+1}}{x_k}.
\tag{9}
\]

\noindent\textbf{[Step 4]}  Substitute (9) into (7) and multiply by $-1$:
\[
\eta\,\inner{\grad f(x_k)}{x_{k+1}-x^{\star}}
= \Breg{x^{\star}}{x_k}-\Breg{x^{\star}}{x_{k+1}}
- \Breg{x_{k+1}}{x_k}.
\tag{10}
\]

\noindent\textbf{[Step 5]}  Use convexity of $f$ (A1) with $y=x^{\star}$:
\[
f(x_k)-f(x^{\star})\le \inner{\grad f(x_k)}{x_k-x^{\star}}.
\tag{11}
\]

\noindent\textbf{[Step 6]}  Add and subtract $\inner{\grad f(x_k)}{x_{k+1}-x^{\star}}$ in (11):
\[
\begin{aligned}
f(x_k)-f(x^{\star})
&\le \inner{\grad f(x_k)}{x_k-x_{k+1}}
   +\inner{\grad f(x_k)}{x_{k+1}-x^{\star}}\\
&= \inner{\grad f(x_k)}{x_k-x_{k+1}}
   +\frac{1}{\eta}\Big(\Breg{x^{\star}}{x_k}
          -\Breg{x^{\star}}{x_{k+1}}-\Breg{x_{k+1}}{x_k}\Big)
   \quad\text{by (10)}.
\end{aligned}
\tag{12}
\]

\noindent\textbf{[Step 7]}  Apply Cauchy–Schwarz and the strong convexity of $h$.
From Lemma~\ref{lem:strong},
\[
\|x_k-x_{k+1}\|^{2}\le 2\,\Breg{x_{k+1}}{x_k}.
\tag{13}
\]
Using $\inner{a}{b}\le\frac{1}{2\eta}\|a\|^{2}+\frac{\eta}{2}\|b\|^{2}$
with $a=\grad f(x_k)$ and $b=x_k-x_{k+1}$,
\[
\inner{\grad f(x_k)}{x_k-x_{k+1}}
\le \frac{1}{2\eta}\|\grad f(x_k)\|_{*}^{2}
   +\frac{\eta}{2}\|x_k-x_{k+1}\|^{2}
\le \frac{1}{2\eta}\|\grad f(x_k)\|_{*}^{2}
   +\eta\,\Breg{x_{k+1}}{x_k}.
\tag{14}
\]

\noindent\textbf{[Step 8]}  Plug (14) into (12):
\[
f(x_k)-f(x^{\star})
\le \frac{1}{2\eta}\|\grad f(x_k)\|_{*}^{2}
   +\eta\,\Breg{x_{k+1}}{x_k}
   +\frac{1}{\eta}\Big(\Breg{x^{\star}}{x_k}
          -\Breg{x^{\star}}{x_{k+1}}-\Breg{x_{k+1}}{x_k}\Big).
\tag{15}
\]

\noindent\textbf{[Step 9]}  The two Bregman terms involving $x_{k+1},x_k$
combine:
\[
\eta\,\Breg{x_{k+1}}{x_k}
-\frac{1}{\eta}\Breg{x_{k+1}}{x_k}
= \Big(\eta-\frac{1}{\eta}\Big)\Breg{x_{k+1}}{x_k}
\le 0,
\]
because $\eta\le 1$ (recall $\eta\le 1/L\le 1$ w.l.o.g.). Hence we can drop this
non‑positive term, obtaining
\[
f(x_k)-f(x^{\star})
\le \frac{1}{2\eta}\|\grad f(x_k)\|_{*}^{2}
   +\frac{1}{\eta}\Big(\Breg{x^{\star}}{x_k}
          -\Breg{x^{\star}}{x_{k+1}}\Big).
\tag{16}
\]

\noindent\textbf{[Step 10]}  Use smoothness (A2) in the form
\[
\|\grad f(x_k)\|_{*}^{2}\le 2L\big(f(x_k)-f(x^{\star})\big),
\tag{17}
\]
which follows from (4) with $y=x^{\star}$ and rearrangement.
Insert (17) into (16):
\[
f(x_k)-f(x^{\star})
\le \frac{L}{\eta}\big(f(x_k)-f(x^{\star})\big)
   +\frac{1}{\eta}\Big(\Breg{x^{\star}}{x_k}
          -\Breg{x^{\star}}{x_{k+1}}\Big).
\tag{18}
\]

\noindent\textbf{[Step 11]}  Choose $\eta\le 1/L$, i.e. $\frac{L}{\eta}\le 1$.
Then move the first term on the right to the left:
\[
\big(1-\frac{L}{\eta}\big)\big(f(x_k)-f(x^{\star})\big)
\le \frac{1}{\eta}\Big(\Breg{x^{\star}}{x_k}
          -\Breg{x^{\star}}{x_{k+1}}\Big).
\tag{19}
\]
Since $1-\frac{L}{\eta}\ge 0$, we obtain
\[
f(x_k)-f(x^{\star})
\le \frac{1}{\eta(1-L\eta)}\Big(\Breg{x^{\star}}{x_k}
          -\Breg{x^{\star}}{x_{k+1}}\Big).
\tag{20}
\]

\noindent\textbf{[Step 12]}  Sum (20) over $k=0,\dots,T-1$:
\[
\sum_{k=0}^{T-1}\big(f(x_k)-f(x^{\star})\big)
\le \frac{1}{\eta(1-L\eta)}\Big(\Breg{x^{\star}}{x_0}
          -\Breg{x^{\star}}{x_T}\Big)
\le \frac{\Breg{x^{\star}}{x_0}}{\eta(1-L\eta)}.
\tag{21}
\]

\noindent\textbf{[Step 13]}  Define the averaged iterate
$\bar{x}_T=\frac{1}{T}\sum_{k=0}^{T-1}x_k$.
By convexity of $f$,
\[
f(\bar{x}_T)-f(x^{\star})
\le \frac{1}{T}\sum_{k=0}^{T-1}\big(f(x_k)-f(x^{\star})\big).
\tag{22}
\]
Combine (21) and (22):
\[
f(\bar{x}_T)-f(x^{\star})
\le \frac{\Breg{x^{\star}}{x_0}}{\eta(1-L\eta)T}.
\tag{23}
\]

\noindent\textbf{[Step 14]}  Setting $\eta=1/L$ (the largest admissible stepsize)
makes $1-L\eta=0$; to avoid division by zero we take the limit $\eta\uparrow 1/L$,
which yields the clean bound
\[
f(\bar{x}_T)-f(x^{\star})\le \frac{L\,\Breg{x^{\star}}{x_0}}{T}.
\tag{24}
\]
This is exactly the statement (2) of Theorem~\ref{thm:md}. ∎

\vspace{1em}
\section*{Rate Derivation (With Constants)}
From (24) we read off the explicit constants:
\[
\boxed{
\text{Rate: }\; f(\bar{x}_T)-f(x^{\star})\le
\frac{L}{T}\,\Breg{x^{\star}}{x_0}.
}
\]
If $h$ is the Euclidean squared norm, $\Breg{x^{\star}}{x_0}=\frac12\|x^{\star}-x_0\|^{2}$,
so the bound becomes
\[
f(\bar{x}_T)-f(x^{\star})\le \frac{L}{2T}\|x^{\star}-x_0\|^{2},
\]
the classic gradient‑descent rate.

\vspace{1em}
\section*{Special Cases}
\begin{itemize}
\item \textbf{Euclidean prox.} $h(w)=\frac12\|w\|^{2}$.
  Mirror Descent reduces to (projected) gradient descent and (24) recovers the
  standard $O(L\|x_0-x^{\star}\|^{2}/T)$ bound.
\item \textbf{Negative entropy.} $h(w)=\sum_{i} w_i\log w_i$ on the simplex.
  $\Breg{x^{\star}}{x_0}=D_{\mathrm{KL}}(x^{\star}\|x_0)$, which can be much smaller
  than the Euclidean diameter, yielding tighter constants for probability
  vectors.
\item \textbf{Adaptive stepsize.} Choosing $\eta_k=1/(Lk)$ leads to the same
  $O(\log T/T)$ bound; the constant stepsize analysis above is simpler and
  optimal for the fixed‑step regime.
\end{itemize}

\end{document}